% QMB_n13_hilbertraum.tex — Kapitel 13: Hilbertraum-Übersetzung
% Inhaltliche Grundlage: Dok. 230 (Mai 2026), Dok. 231; eigene Darstellung
\chapter{Zwei Sprachen, ein Inhalt: FFGFT und Hilbertraum}
\label{ch:hilbert}

\section{Die Behauptung}

Die FFGFT beschreibt Quantensysteme in der Sprache von Moden auf einem
Torus. Die Standard-Quantenmechanik beschreibt sie in der Sprache von
Vektoren in einem Hilbertraum. Dieses Kapitel zeigt, dass beide Sprachen
denselben Inhalt tragen: Jede Rechnung, die in der einen möglich ist,
ist in der anderen möglich, mit demselben Ergebnis. Die Übersetzung ist
keine Näherung und keine Spezialisierung, sondern eine bijektive
Zuordnung von Zuständen, Operatoren, Gattern, Produkträumen und
Messvorschriften.\status{B}

Das ist methodisch die wichtigste Aussage dieses Buches. Sie bedeutet,
dass die FFGFT die Quantenmechanik nicht ersetzt, sondern erweitert:
Beide Formalismen liefern bis auf eine $\xi$-Korrektur der Größenordnung
$10^{-5}$ dieselben messbaren Vorhersagen. Sie unterscheiden sich in der
Deutung einiger Grundbegriffe --- Born-Regel, Nichtlokalität, der Status
von $\hbar$, die Rolle der Raumzeit, Singularitäten. Wer in
Hilbertraum-Sprache zu Hause ist, kann die FFGFT benutzen, ohne seine
Werkzeuge zu wechseln. Die Brücke ist die Übersetzungstabelle.

\section{Wo wir uns geometrisch befinden}

Bevor die Tabellen kommen, ein Bild der Geometrie, in der die Übersetzung
stattfindet. Wir werden mit Zylinderkoordinaten rechnen, und der Zylinder
ist nicht die volle FFGFT-Geometrie, sondern eine bequeme Reduktion.

\paragraph{Die volle Geometrie ist vierdimensional.} Der Raum, auf dem
alle Aussagen der FFGFT leben, ist der Torus $T^4 = T^3\times S^1$: drei
räumliche Wicklungsrichtungen und eine zeitliche. Die Wellenfunktion
einer Mode $\varphi_{n,l,j}$ ist eine Funktion auf diesem Torus. Aus
seiner Geometrie folgen die Massen, die Feinstrukturkonstante, die
kosmologischen Größen --- alles, was die FFGFT von der reinen
Quantenmechanik unterscheidet.

\paragraph{Ein einzelnes Qubit braucht weniger.} Ein Zwei-Niveau-System
zu fester Zeit, ohne Massendynamik, aktiviert die meisten dieser
Wicklungsrichtungen nicht. Zwei Drehfreiheitsgrade bleiben: einer, der
kodiert, wo zwischen $|0\rangle$ und $|1\rangle$ wir uns befinden (das
wird $z$), und einer für die Phase (das wird $\theta$). Der natürliche
Raum dafür ist ein 2-Torus $T^2 = S^1\times S^1$.

\paragraph{Vom 2-Torus zum Zylinder.} Ist der Hauptradius $R$ des
2-Torus groß gegen den Schlauchradius, sieht die Geometrie lokal wie ein
Zylinder $\mathbb{R}\times S^1$ aus. Eine der beiden Schleifen ist
aufgeschnitten: Aus zwei Kreisen wird eine Gerade und ein Kreis. In
Zylinderkoordinaten $(z,r,\theta)$ ist $z$ deshalb eine lineare Achse
von $-1$ bis $+1$, nicht mehr zyklisch.

\paragraph{Vom Zylinder zur Bloch-Sphäre.} Zieht man den Zylinder oben
und unten zu je einem Punkt zusammen, entsteht eine Sphäre $S^2$: die
Bloch-Sphäre der Standard-Quantenmechanik. Sie ist die zweite,
weitergehende Reduktion.

\begin{center}
\small\renewcommand{\arraystretch}{1.3}
\begin{tabular}{p{2.9cm}cp{3.6cm}}
\toprule
\textbf{Geometrie} & \textbf{Dim.} & \textbf{Verwendung} \\
\midrule
4-Torus $T^4$ & 4 & volle FFGFT: Massen, $\alpha$, Kosmologie \\
3-Torus $T^3$ & 3 & Modenraum bei fester Zeit \\
2-Torus $T^2$ & 2 & Einzelqubit ohne Reduktion \\
Zylinder $\mathbb{R}\times S^1$ & 2 & Limes $R\to\infty$, lokale Näherung \\
Bloch-Sphäre $S^2$ & 2 & Pole zusammengezogen; Standard-QM \\
\bottomrule
\end{tabular}
\end{center}

\paragraph{Was die Reduktionen kosten.} Jede Stufe verliert etwas. Vom
4-Torus zum 3-Torus geht die Zeitwicklung verloren und damit jede
Aussage über Massen. Vom 3-Torus zum 2-Torus geht eine räumliche
Wicklung verloren und damit die Verbindung zur dritten Leptongeneration
und zu vielen Spinstrukturen. Vom 2-Torus zum Zylinder geht eine
kompakte Schleife verloren --- das ist die schwächste Reduktion, und
genau hier sitzt der fraktale Korrekturfaktor
$K_{\mathrm{frak}} = 1-100\xi\approx0{,}9867$, der in den
Bit-Flip-Drehungen auftaucht. $K_{\mathrm{frak}}$ ist die Erinnerung des
Zylinders an seine Torusherkunft.

Die Übersetzungstabellen dieses Kapitels gelten für den Qubit-Sektor,
also für die unteren drei Zeilen. Die Standard-QM arbeitet auf der
Bloch-Sphäre, die FFGFT auf dem Zylinder oder dem 2-Torus. Die
Übersetzung zwischen diesen Stufen ist verlustfrei, sofern man die
kleinen geometrischen Korrekturen ($K_{\mathrm{frak}}$, $\Delta$CHSH)
mitführt. Der 4-Torus liegt eine Stufe höher und liefert, was der
Hilbertraum nicht liefern kann: die Werte der Eingabeparameter.

\begin{laierspur}
Man kann eine Landkarte auf verschiedene Weisen zeichnen: als Globus,
als abgerollten Zylinder, als flaches Blatt. Alle zeigen dieselben
Städte an denselben Stellen. Aber nur der Globus zeigt, dass man von
Osten aus wieder im Westen ankommt. Die Bloch-Sphäre der
Quantenmechanik ist das flache Blatt; der Torus der FFGFT ist der Globus.
Für die Navigation zwischen zwei Städten reicht das Blatt. Für die
Frage, warum die Welt so groß ist, wie sie ist, braucht man den Globus.
\end{laierspur}

\section{Zustände}

In der Standard-QM ist ein Qubit ein Vektor in $\mathbb{C}^2$:
\begin{equation}
  |\psi\rangle = \alpha|0\rangle + \beta|1\rangle,\qquad
  |\alpha|^2+|\beta|^2 = 1,\qquad \alpha,\beta\in\mathbb{C}.
\end{equation}
Mit globaler Phasenfreiheit ist das ein Punkt auf der Bloch-Sphäre
(Hopf-Faserung $S^3\to S^2$).

In der FFGFT ist dasselbe Qubit ein Punkt $(z,r,\theta)$ auf dem
Zylinder:
\begin{equation}
  z\in[-1,1],\quad r\in[0,1],\quad \theta\in[0,2\pi),\qquad z^2+r^2 = 1.
\end{equation}
$z$ ist die Projektion auf die Berechnungsachse, $r$ die radiale
Superpositionsamplitude, $\theta$ die Phase.

Die Brücke zwischen beiden Darstellungen:
\begin{equation}
  \boxed{\;
  \alpha = \sqrt{\tfrac{1+z}{2}}\;e^{\,i\theta/2},\qquad
  \beta = \sqrt{\tfrac{1-z}{2}}\;e^{-i\theta/2}.\;}
  \label{eq:state-bridge}
\end{equation}
Umgekehrt:
\begin{equation}
  z = |\alpha|^2-|\beta|^2,\qquad
  r = 2|\alpha\beta|,\qquad
  \theta = \arg\alpha-\arg\beta.
\end{equation}
Die Konsistenz prüft man in vier Zeilen: Die Normierung
$|\alpha|^2+|\beta|^2 = \frac{1+z}{2}+\frac{1-z}{2} = 1$ stimmt; die
Zylinderbedingung $z^2+r^2 = (|\alpha|^2-|\beta|^2)^2+4|\alpha|^2|\beta|^2 = 1$
stimmt; $|0\rangle$ entspricht $(1,0,\ast)$ und $|1\rangle$ entspricht
$(-1,0,\ast)$, mit beliebiger Phase an den Polen.

Der begriffliche Unterschied liegt nicht in den Zahlen. In der QM ist
$|\alpha|^2$ eine Wahrscheinlichkeit. In der FFGFT ist $r^2 = 1-z^2$
eine Energiedichte der radialen Konfiguration. Beide Größen sind bei
großer Statistik numerisch gleich; konzeptuell ist die FFGFT ein
deterministisches geometrisches Modell, dessen statistische Vorhersagen
mit der Born-Regel zusammenfallen.

\section{Operatoren}

Die drei Pauli-Matrizen erzeugen alle Ein-Qubit-Operationen:
\begin{equation}
  \sigma_x = \begin{pmatrix}0&1\\1&0\end{pmatrix},\quad
  \sigma_y = \begin{pmatrix}0&-i\\i&0\end{pmatrix},\quad
  \sigma_z = \begin{pmatrix}1&0\\0&-1\end{pmatrix},
\end{equation}
mit der Algebra $\sigma_i\sigma_j = \delta_{ij}\mathbb{1}+i\epsilon_{ijk}\sigma_k$.

Auf dem Zylinder entsprechen ihnen drei geometrische Drehungen:
$\sigma_z$ ist die Rotation um die Zylinderachse, also die Phasenrotation
$\theta\to\theta+\pi$; $\sigma_x$ ist die Drehung in der $(z,r)$-Ebene um
den Winkel $\pi$, der Bit-Flip; $\sigma_y$ ist die dazu orthogonale
Drehung, die Bit-Flip und Phasenrotation kombiniert.

\begin{center}
\small\renewcommand{\arraystretch}{1.5}
\begin{tabular}{cp{2.6cm}p{3.6cm}}
\toprule
\textbf{Operator} & \textbf{Hilbertraum} & \textbf{FFGFT-Geometrie} \\
\midrule
$\sigma_z$ & $\mathrm{diag}(1,-1)$ & Phasenrotation $\theta\to\theta+\pi$ \\
$\sigma_x$ & $\begin{pmatrix}0&1\\1&0\end{pmatrix}$ & $(z,r)$-Drehung um $\pi K_{\mathrm{frak}}$ \\
$\sigma_y$ & $\begin{pmatrix}0&-i\\i&0\end{pmatrix}$ & Drehung mit Phasenkippung \\
$H$ & $\frac{1}{\sqrt2}\begin{pmatrix}1&1\\1&-1\end{pmatrix}$ & $z\leftrightarrow r$ vertauschen, $\theta\to\theta+\pi/2$ \\
$T$ & $\mathrm{diag}(1,e^{i\pi/4})$ & Phasenrotation $\theta\to\theta+\pi/4$ \\
\bottomrule
\end{tabular}
\end{center}

Die einzige FFGFT-spezifische Abweichung ist der Faktor
$K_{\mathrm{frak}}\approx0{,}9867$ in der $\sigma_x$-Drehung. Er folgt aus
der fraktalen Raumzeitdimension $D_f = 2{,}999867$ und ist eine prüfbare
Vorhersage: Alle $\pi$-Gatter weichen systematisch um
$\Delta\alpha\approx0{,}013\,\%$ vom idealen Wert ab.\status{S} Im
Hilbertraum-Formalismus wäre das eine Ad-hoc-Korrektur, die als
Hardware-Rauschen eingebaut werden müsste; in der FFGFT folgt sie aus
der Geometrie.

\section{Gatter}

Jedes unitäre Ein-Qubit-Gatter $U\in U(2)$ ist eine Bloch-Drehung:
\begin{equation}
  U = e^{i\alpha}R_{\hat n}(\phi)
    = e^{i\alpha}\bigl(\cos\tfrac{\phi}{2}\,\mathbb{1}
      - i\sin\tfrac{\phi}{2}\,\hat n\cdot\vec\sigma\bigr).
\end{equation}
In der FFGFT ist das die Rotation des Punktes $(z,r,\theta)$ um die Achse
$\hat n$ um den Winkel $\phi$. Die Übersetzung ist exakt.

Das CNOT-Gatter
\begin{equation}
  \mathrm{CNOT} =
  \begin{pmatrix}1&0&0&0\\0&1&0&0\\0&0&0&1\\0&0&1&0\end{pmatrix}
\end{equation}
übersetzt sich als Vorschrift: Wenn $z_A = -1$ (Kontrollqubit auf
$|1\rangle$), führe den Bit-Flip auf Qubit B aus. Das ist eine geometrische
Konditionierung --- die Drehung des einen Zylinderpunkts hängt vom
$z$-Wert des anderen ab. Topologisch ist das eine Verkettung der beiden
Zylinder, und das ist die Verschränkung.

Die Menge $\{H,T,\mathrm{CNOT}\}$ ist nach dem Solovay-Kitaev-Theorem
universell für die Quantenmechanik.\status{E} Da jedes dieser Gatter in der
FFGFT als geometrische Operation darstellbar ist, ist die FFGFT
vollständig universell für Quantenberechnung: Es gibt keine QM-Aussage,
die nicht in FFGFT-Sprache formuliert werden kann.\status{B}

\section{Produkträume und Verschränkung}

Für $n$ Qubits ist der Hilbertraum
$\mathcal{H}_n = \bigotimes_{i=1}^n\mathbb{C}^2$ mit Dimension $2^n$,
und ein allgemeiner Zustand lautet
\begin{equation}
  |\Psi\rangle = \sum_{x\in\{0,1\}^n}\alpha_x|x\rangle,\qquad
  \sum_x|\alpha_x|^2 = 1.
\end{equation}
In der FFGFT entsprechen $n$ Qubits $n$ Modenkonfigurationen
$(z_i,r_i,\theta_i)$ auf einer gemeinsamen Torusgeometrie. Verschränkung
ist eine topologische Verbindung: Verschränkte Qubits sind Knoten der
Modenkonfiguration, die nicht auf Produktform reduziert werden können.
Der Bell-Zustand $|\Phi^+\rangle = \frac{1}{\sqrt2}(|00\rangle+|11\rangle)$
ist eine topologisch verbundene Zylinderkonfiguration, in der die
einzelnen $(z_i,r_i,\theta_i)$ nicht unabhängig sind.

Beide Formalismen reproduzieren das exponentielle Wachstum $\dim = 2^n$:
im Hilbertraum durch das Tensorprodukt, in der FFGFT durch $n$
unabhängige Modenkonfigurationen mit topologischen Verbindungen. Die
Zahl der reellen Parameter pro Zustand ist in beiden Fällen $2\cdot2^n-2$.

Die CHSH-Vorhersage: Die Standard-QM gibt $|\langle\mathrm{CHSH}\rangle|
\le2\sqrt2$. Die FFGFT sagt eine kleine, geometrisch motivierte Abweichung
voraus, $\Delta\mathrm{CHSH}\approx\xi/(2\pi)\approx10^{-5}$ pro
Messung.\status{S} Sie liegt unter der aktuellen Hardware-Rauschschwelle;
die Unterscheidung von der kumulativen Größe aus Dok.~022 ist in
Kapitel~\ref{ch:chsh_exp} ausgeführt.

\section{Messung}

Eine Messung in der Berechnungsbasis liefert in der QM den Wert $0$ mit
Wahrscheinlichkeit $|\alpha|^2$ und den Wert $1$ mit Wahrscheinlichkeit
$|\beta|^2$; der Zustand kollabiert auf $|0\rangle$ oder $|1\rangle$.

In der FFGFT ist Messung kein probabilistischer Kollaps, sondern eine
geometrische Wechselwirkung, die den Punkt $(z,r,\theta)$ auf das
nächstgelegene stabile Extrem $z = \pm1$ treibt. Das ist eine
deterministische Bewegung im Modenraum --- wegen der unbekannten Phase
$\theta$ und kleinster Variationen der Anfangsbedingungen statistisch
aber ununterscheidbar vom Verhalten der Born-Regel. Bei großer Statistik
konvergiert die Messverteilung:
\begin{equation}
  P_{\mathrm{FFGFT}}(z = +1) = \tfrac{1+z}{2} = |\alpha|^2 = P_{\mathrm{QM}}(0).
\end{equation}
Die Formalismen liefern identische statistische Vorhersagen.\status{B}
Der Streit über lokalen Realismus löst sich in der FFGFT geometrisch,
weil die Verschränkung als topologische Verbindung und nicht als
nichtlokale Korrelation dargestellt wird (Kapitel~\ref{ch:bell_ffgft}).

\section{Zeitentwicklung}

Die Zeitentwicklung der QM ist die Schrödinger-Gleichung
$i\hbar\,\partial_t|\psi\rangle = H|\psi\rangle$. In der FFGFT ist sie die
deterministische Bewegung der Modenkonfiguration $\varphi_{n,l,j}$ auf
$T^3\times\mathbb{R}$. Die Modengleichung reduziert sich im
Niederenergielimes $E\ll E_\xi$ auf die Schrödinger-Gleichung:
\begin{equation}
  i\hbar\,\frac{\partial\varphi_{n,l,j}}{\partial t}
  \;\xrightarrow{\;E\ll E_\xi\;}\;
  -\frac{\hbar^2}{2m}\nabla^2\varphi_{n,l,j}+V\varphi_{n,l,j}.
\end{equation}
Bei höheren Energien treten $\xi$-Korrekturen auf, die im
Hilbertraum-Formalismus als Modifikationen des Hamiltonoperators
erscheinen.

\section{Was übersetzbar ist und was nicht}

Alle Hilbertraum-Aussagen sind in FFGFT abbildbar: Zustände, Operatoren
und Gatter (mit oder ohne $K_{\mathrm{frak}}$), Tensorprodukte,
Verschränkung, Bell- und GHZ-Zustände, Messung, unitäre Evolution und
alle Algorithmen. Umgekehrt entspricht jede $(z,r,\theta)$-Konfiguration
eindeutig einem $(\alpha,\beta)$-Vektor, jede geometrische Drehung einem
unitären Operator, jede topologische Verbindung einem verschränkten
Vektor. Jede Berechnung, die in einem Formalismus durchführbar ist, ist
im anderen durchführbar --- mit identischem Ergebnis.

Die Übersetzbarkeit gilt für Aussagen über Quantensysteme. Sie gilt
nicht für Aussagen über die Eingabewerte der Quantenmechanik:
\begin{itemize}
  \item Der Parameter $\xi = 4/30000$ ist im Hilbertraum ein externer
        Eingabewert. In der FFGFT folgt er aus der Geometrie
        (Kapitel~\ref{ch:grundlage}).
  \item Die Fermionmassen erscheinen im Hilbertraum als Yukawa-Kopplungen.
        In der FFGFT folgen sie aus $m_i = r_i\,\xi^{p_i}\,v$.
  \item Die Feinstrukturkonstante ist in der FFGFT ableitbar, im
        Hilbertraum eingegeben.
  \item $K_{\mathrm{frak}}$ erscheint in FFGFT-Gattern als
        Geometrieabweichung; im Hilbertraum müsste er als Rauschen
        modelliert werden.
\end{itemize}
Das ist die strukturelle Asymmetrie: Der Hilbertraum hat alle Werkzeuge
für QM-Berechnungen, aber keine Erklärung für die Eingabewerte. Die FFGFT
liefert die Werkzeuge und die Eingabewerte aus einer einzigen Geometrie.
In diesem Sinn ist die Quantenmechanik eine Untermenge der FFGFT: die
FFGFT, eingeschränkt auf die Frage, wie sich Systeme verhalten, ohne die
Frage, warum die Konstanten diese Werte haben.

Die FFGFT erhebt keinen Anspruch, die einzige korrekte Darstellung zu
sein. Sie ist eine Darstellung, die zusätzlich die Werte der
fundamentalen Konstanten liefert. Wer diesen Mehrwert nicht braucht,
kann sie ignorieren; die Standard-QM bleibt gültig.

\section{Operationale Äquivalenz, interpretative Divergenz}

Die vollständige Übersetzbarkeit darf nicht als ontologische Identität
gelesen werden. Operational stimmen FFGFT und Standard-QM in jedem
messbaren Ergebnis bis auf $\xi$-Korrekturen überein. Interpretativ
weichen sie an fünf Stellen ab, an denen die Standardlesart Aussagen
trifft, die in der FFGFT als Artefakte der gewählten Trägerannahmen
gelesen werden:

\begin{enumerate}
  \item \textbf{Born-Regel als irreduzibler Zufall.} Standardlesart:
        $|\psi|^2$ ist ontologische Wahrscheinlichkeit, der Messausgang
        fundamental zufällig. FFGFT: $|\psi|^2$ ist die korrekte Verteilung
        der Beobachtungen, entsteht aber aus einer deterministischen
        Polübergangsdynamik. Die Stochastizität sitzt beim Messapparat,
        nicht im Naturprozess.
  \item \textbf{Nichtlokalität als Eigenschaft der Natur.} Standardlesart
        der Bell-Korrelationen: Fernwirkung im $\mathbb{R}^3$. FFGFT:
        topologische Verbindung auf $T^4$, geometrische Nähe in einer
        geschlossenen Geometrie.
  \item \textbf{$\hbar$ als Primitive.} Standardlesart: fundamentale
        Naturkonstante. FFGFT: SI-Konversionsfaktor zwischen Energie- und
        Zeitmaß, der aus der Abschlussbedingung $S = h$ auf $T^4$ folgt.
        Numerisch identisch, ontologisch abgeleitet.
  \item \textbf{Raumzeit als Bühne.} Standardlesart der Relativitäts-
        theorie: gegebene Mannigfaltigkeit. FFGFT: aus der
        $T^4$-Identifikationsgeometrie abgeleitet. Fragen wie
        \enquote{was war vor dem Urknall} sind Artefakte der
        $\mathbb{R}^3\!+\!t$-Annahme.
  \item \textbf{Singularitäten als Realitäten.} Standardlesart: Schwarzes
        Loch und Urknall als echte Singularitäten. FFGFT: Da $T^4$
        geschlossen ist, sind sie Endpunkte des Trägermodells, nicht der
        Physik. Außerhalb der Singularitätsregion stimmen die Vorhersagen
        überein.
\end{enumerate}

Diese fünf Punkte sind keine empirischen Behauptungen gegen die
Standard-QM. Sie sind Interpretationsdifferenzen, die sichtbar werden,
sobald nach dem ontologischen Träger gefragt wird. Wer nur mit den Zahlen
rechnet, kann sie ignorieren. Wer wissen will, warum die Zahlen gerade
diese sind, für den werden sie zu prüfbaren strukturellen Aussagen.

Innerhalb der heutigen Messpräzision --- NISQ-Niveau $\sim10^{-2}$,
optische Präzisionsexperimente $\sim10^{-12}$ bei Bell-Tests --- liefern
beide dieselben Vorhersagen. Die $\xi$-Korrektur $\sim10^{-5}$ ist mit
gegenwärtiger Hardware nicht auflösbar; das ist auf IBM Heron~r2
empirisch geprüft (Kapitel~\ref{ch:chsh_exp}).

\section{Feldtheoretische Erweiterung}

Die Übersetzung lässt sich vom Qubit-Sektor auf Quantenfeldzustände
ausdehnen. Der Fock-Raum der Standard-QFT --- die direkte Summe aller
$n$-Teilchen-Räume --- entspricht dem Modenspektrum des Trägers, wobei
die Besetzungszahlen als Wicklungszahlen erscheinen. Erzeugungs- und
Vernichtungsoperatoren sind Übergänge zwischen benachbarten
Wicklungszuständen derselben Mode; Kommutatoren folgen aus der
Periodizität. Das erzeugt keine neuen Vorhersagen im perturbativen
Bereich, strukturiert aber den Aufbau des Standardmodells geometrisch um:
Die Teilchensorten sind Modenklassen des Gitters, nicht unabhängige
Felder mit je eigener Kopplung.\status{S}

\begin{keypoint}[Was dieses Kapitel festhält]
Zustände, Operatoren, Gatter, Verschränkung, Messung und Zeitentwicklung
sind zwischen FFGFT und Hilbertraum bijektiv übersetzbar. Die einzige
formale Abweichung ist $K_{\mathrm{frak}}$ in den $\pi$-Gattern; die
einzige statistische ist $\Delta$CHSH. Der Mehrwert der FFGFT liegt nicht
in anderen Rechenergebnissen, sondern darin, dass die Eingabewerte der
Quantenmechanik aus derselben Geometrie folgen.
\end{keypoint}

\quelle{Dok.~230 (Übersetzbarkeit, Mai 2026), Dok.~231 (feldtheoretische Erweiterung), Dok.~147 §1.2, Dok.~175 §3--4, Dok.~202 §22}
